\section{Independent proof of compatibility of the two reductions}
\label{sec:independent-composition}

The exact combination is mathematically justified; the fact that its
edited rows have disjoint labels is only one part of the justification.
We independently check the rows used by every exceptional block, all
incoming ports needed by the structural invariants, the complete edit
sets, and equality of the two possible orders of construction.

Let $S$ denote the repaired boundary240 table, let $B$ denote the
entry-context reduction deleting state $218$, and let $D$ denote the
delayed-erasure reduction deleting state $169$.  All labels in this
section are inherited labels: start $95$, dispatcher $96$, and halt
$254$.  The delivered table $N$ satisfies the literal identity
\[
             N=D(B(S))=B(D(S)).
\]
The two intermediate tables $B(S)$ and $D(S)$ each have $239$ working
states; $N$ has $238$ working states and $476$ transition rows.
This identity is equality of complete tables, not merely equivalence
of sampled executions.  The independent audit also verifies that
all $238$ working states are reachable in the directed transition
graph from $95$.  Graph reachability is used for table accounting;
it makes no assertion about which states occur in the blank run.

\subsection{The delayed-erasure edge remains valid}

Write $T$ for the complete binary tape and $h$ for the head coordinate.
In the predecessor $B(S)$ the following are structural invariants:
\[
\begin{array}{c|l}
q&\text{required tape facts}\\\hline
67&T(h-1)=0\\
80&T(h-2)=0,\quad T(h-1)=1\\
91&T(h-3)=0,\quad T(h-2)=T(h-1)=1\\
169&T(h-1)=0.
\end{array}
\]
They hold on every binary tape started in state $95$.  The only
incoming rows to $67$ are $52/0$ and $124/0$; both write $0$ and move
right.  The only incoming rows to $80$, $91$, and $169$ are
$67/1$, $80/1$, and $147/0$, respectively.  These facts have been
rechecked against the complete boundary239 table.  The boundary
transformation leaves every one of these rows unchanged.

The additional selector fact
\begin{equation}
q=169\ \land\ T(h)=1\quad\Longrightarrow\quad T(h+1)=1
\label{eq:new-selector-guard}
\end{equation}
requires the interpreter's representation invariant.  The old guard
certificate for the unmodified rotation240 source cannot simply be
reused: the boundary construction changes source paths.  The new
guard replay therefore checks the exact boundary239 table against
the exact changed source, including the literal blank initialization
and every unbounded instruction case.  No execution rule is added
to the symbolic engine.  Every selector entry must expose a literal
run or the known zero exterior; every entry reading $1$ must exhibit
a uniformly positive run of length at least two.  The generic
finite-list block lemmas are rederived and checked not to hide a
selector visit.  State $169$ has no self-loop, so unary-loop
acceleration cannot hide a selector visit either.

The fresh run passes all $8,386$ guarded instruction cases using
$2,247,513$ symbolic steps.  It finds $96$ symbolic selector entries,
of which $32$ read $1$.  At each of those $32$ entries the right-hand
run has length $G+K$, with $G\geq0$ the additional header gap and
$23\leq K\leq59$; all are increment-register-$4$ sites.  The exact
$409,691$-step boundary239 startup has $18$ selector entries, including
six reading $1$, and every one satisfies the required implication.
These are fresh measurements for the changed source, rather than
the predecessor's $21\leq K\leq59$ interval.  The receipt is
\texttt{boundary239\_selector\_guard.json}.

With this guard, put $\phi(169)=80$ and $\phi(q)=q$ otherwise.
Corresponding configurations retain the exact same head and entire
tape.  All old transitions outside $80$, $91$, and $169$ agree in
one step under $\phi$.  The complete exceptional identities are:
\[
\begin{array}{c|l|r|r|c}
\text{old row}&\text{tape assumptions}&\text{old steps}&\text{new steps}
&\text{endpoint}\\\hline
80/0&T(h-1)=1&1&5&(25,h+1)\\
80/1&\text{none}&1&1&(91,h+1)\\
169/0&T(h-1)=0&1&3&(57,h+1)\\
169/1&T(h-1)=0,\ T(h+1)=1&1&5&(56,h+1)\\
91/0&\text{none}&1&1&(106,h+1)\\
91/1&T[h-3,h)=011&3&11&(2,h-3).
\end{array}
\]
The read bit specified in the first column is always part of the
assumptions.  Every identity equates the complete final tapes.
Unmentioned cells are independent free variables, and the checker
rejects any attempt to read one.  In particular, these are exact
local identities valid for arbitrary exterior tapes.

All rows executed in these six old/new comparisons are unchanged
by $B$.  The independent checker replays the identities both from
$S$ to $D(S)$ and from $B(S)$ to $N$, and obtains identical traces.
Thus the only source-sensitive premise was the selector guard,
which is discharged afresh on $B(S)$.

\subsection{The boundary edge also survives delayed erasure}

The $B$ reduction uses a different part of control.  At old $179$,
the bit $\ell=T(h-1)$ is always $0$: its only incoming rows are
$161/0$, $236/0$, and $239/0$, all writing $0$ and moving right.
Old $218$ receives tag $1$ from $187/1$ and $220/1$, or tag $0$
from $237/0$.  Old $223$ can be entered only from $218/0$ and is
kept entirely inside the corresponding three-step comparison block.
At checkpoints the old/new control relation is
\[
179,\ell=0\mapsto179,\qquad
218,\ell=0\mapsto223,\qquad
218,\ell=1\mapsto179,
\]
with all other checkpoint controls unchanged and old $223$
excluded from the checkpoints.  The tape and head are identical.

The six comparison identities below read only the cells $h-1,h$,
leave the entire tape unchanged, and finish at $h+1$:
\[
\begin{array}{r|r|r|r|r|r|r}
\text{old state}&\ell&\text{read}&\text{new state}&
\text{old steps}&\text{new steps}&\text{endpoint}\\\hline
179&0&0&179&1&1&105\\
179&0&1&179&1&5&105\\
218&0&0&223&3&1&23\\
218&0&1&223&1&1&70\\
218&1&0&179&3&1&105\\
218&1&1&179&1&3&70.
\end{array}
\]
Every required incoming port and executed row survives $D$ unchanged.
The independent checker replays this proof both from $S$ to $B(S)$
and from $D(S)$ to $N$.  This boundary equivalence holds for every
initial binary tape at start $95$; it requires no semantic selector
guard and no restriction on source-register values.

\subsection{Why positive finite blocks prove halting equivalence}

Suppose corresponding configurations satisfy the appropriate
checkpoint invariant.  At an ordinary checkpoint one step on
each side restores the relation.  At an exceptional checkpoint
one of the displayed positive finite blocks restores it with
equal complete tapes and head positions.  None of the blocks
passes through the halt state before its endpoint.  Induction
therefore constructs corresponding checkpoints for as many
blocks as either deterministic computation performs.  Positivity
prevents an infinite execution on one side from being represented
by only finitely many steps on the other.  A finite halting
execution corresponds to a finite halting execution, and an
infinite execution corresponds to an infinite execution.

Applied to the freshly guarded delayed-erasure edge, this proves
\[
             N\downarrow\quad\Longleftrightarrow\quad B(S)\downarrow
\]
from the blank tape.  The independent boundary edge gives
$B(S)\downarrow\leftrightarrow S\downarrow$.  The separate source
equivalence, decoder/runtime, and startup arguments connect $S$
to the calibrated rotation240 source.  Composition yields the
requested consistency equivalence for the literal $238$-state table.

The finite block arguments and natural-number/list inductions are
ordinary mathematical proofs formalizable in ZFC.  Their executable
certificates are not a newly compiled Lean endpoint for $N$.
The existing external Lean endpoint for rotation240 identifies its
nonhalting with absence of a finite refutation in the specified
recursive ZFC proof calculus.  The new composition preserves that
halting predicate.  A proof of the entire equivalence internal to
PA, and an external Lean compilation for this exact new table,
remain distinct artifacts and are not asserted here.
